\rednote{Stand nach VL vom Freitag 13.11.2026.}

\subsection{Grundlegende Definitionen}

\begin{definition}
  Eine \begriff{Ring} (mit Eins) ist eine Menge \(R\) zusammen mit zwei Verknüpfungen \(+\) und \(\cdot\) sodass gilt:
  \begin{enumerate}[(i)]
  \item \((R,+)\) ist abelsche Gruppe.
  \item \((R,\cdot)\) ist Monoid.
  \item Es gelten die folgenden beiden Distributivgesetze:
    \begin{align*}
      &\forall a,b,c\in R.~ (a+b)\cdot c =a\cdot c+b\cdot c \\
      &\forall a,b,c\in R.~ a\cdot (b+c)=a\cdot b+a\cdot c
    \end{align*}
  \end{enumerate}
  Ein Ring heißt \emph{kommutativ}, wenn \(\cdot\) kommutativ ist.
  Wir gehen stets davon aus, dass das der Fall ist und erwähnen Ausnahmen explizit.
\end{definition}

\begin{konvention}
  Das Neutralelement bezüglich \(+\) bezeichnen wir stets mit \(0\) und das Neutralelement bezüglich \(\cdot\) mit \(1\). Wie gehabt schreiben wir \(-a\) für das additiv Inverse zu \(a\in R\) und \(a^{-1}\) für das multiplikativ Inverse. Wir lassen \(\cdot\) wie gehabt manchmal weg und verwenden die Regel ``Punkt vor Strich'', dass also \(\cdot\) stärker bindet als \(+\), so ist zum Beispiel:
  \[
    ac+b=(a\cdot c)+b
  \]  
\end{konvention}

\begin{beispiel}
  \begin{enumerate}[(a)]
  \item \(\{0\}\) ist der sogenannte \begriff{Nullring}, der einzige Ring in dem \(1=0\) gilt.
  \item \((\Z,+,\cdot)\), \((\Q,+,\cdot)\), \((\R,+,\cdot)\), \((\C,+,\cdot)\) sind Ringe.
  \item \(M_n(R)\coloneq R^{n\times n}\) ist nicht-kommutativer Ring mit Matrixmultiplikation und komponentenweiser Addition.
  \item Die Potenzmenge \(\mathcal P(M)\) einer beliebigen Menge \(M\) ist ein Ring mit Durchschnitt als Multiplikation und symmetrischer Differenz als Addition.
  \item Wie bei den Gruppen ist das Produkt \(R\times S\) zweier Ringe \(R,S\) mit komponentenweiser Verknüpfung wieder ein Ring:
  \begin{align*}
  & (r,s)+(r^\prime,s^\prime)\coloneq (r+r^\prime,s+s^\prime) \\
  & (r,s)\cdot (r^\prime,s^\prime)\coloneq (r\cdot r^\prime,s\cdot s^\prime)\rlap{.}
\end{align*}
\(R\times S\) ist genau dann kommutative, wenn \(R\) und \(S\) kommutativ sind.
Allgemeiner ist für eine Menge \(M\) und einen Ring \(R\) auch \(M\to R\) mit komponentenweisen/punktweisen Verknüpfungen wieder ein Ring und \(\prod_{i\in I}R_i\) für eine Familie \((R_i)_{i\in I}\) von Ringen.
  \end{enumerate}
\end{beispiel}

\begin{bemerkung}
  Sei \(R\) ein Ring, dann gilt: \(-a=(-1)\cdot a\).
\end{bemerkung}

\begin{beweis}
  \(0=(1+(-1))\cdot a=a+(-1)\cdot a \Rightarrow -a=(-1)\cdot a\).
\end{beweis}

\begin{definition}
  Sei \(R\) Ring.
  Die \begriff{Einheitengruppe} \(R^\times\)\index{\(R^\times\)} ist die Gruppe der invertierbaren Elemente des Monoids \((R,\cdot)\).
\end{definition}

\begin{definition}
  Sei \(R\) ein Ring.
  \begin{enumerate}
  \item Ein Element \(r\in R\) heißt \begriff{Nullteiler}, wenn \(s\in R\) existiert, sodass \(r\cdot s=0\) gilt.
  \item \(R\) heißt \begriff{nullteilerfrei}, wenn für jeden Nullteiler \(r\in R\) bereits \(r=0\) gilt.
  \end{enumerate}
\end{definition}

\begin{beispiel}
  \(\Z\) ist nullteilerfrei, \(\Z\times \Z\) hat zum Beispiel den Nullteiler \((1,0)\).
\end{beispiel}



\begin{definition}
Sei \(R\) ein Ring. Die zugrundeliegende Menge des \begriff{Polynomring}s\index{\(R[X]\)} ist gegeben durch:
\[
  R[X]\coloneq \{a\in R^\N\mid \exists n\in \N. \forall k\geq n. a_k=0\}
\]
Die Addition ist komponentenweise gegeben, also für \(a,b\in R[X]\) sei:
\[
  (a+b)_i\coloneq a_i+b_i
\]
Die Multiplikation ist gegeben durch die Cauchy-Produktformel:
\[
  (a\cdot b)_i\coloneq \sum_{k,l\in\N, k+l=i}a_k\cdot b_l
\]
\end{definition}
